\section{LLM-Assisted Assembly Transformation}\label{sec:llm}
\subsection{What the Model Is Asked to Change}
In this mode the LLM receives assembly and proposes a new instruction
sequence intended to preserve behavior. It supports the EQ, NOP, BCF,
and IBP families introduced in Section~\ref{sec:background}; data
classification instead returns only locations. PII and wrapping remain
native operations because their layout and encoding constraints are not
implemented in the LLM interface's requests and validators. This is an
implementation boundary, not an intrinsic limitation of language models.

Zone prepares an assembly listing with identified data retained separately.
The model proposes instruction changes while deterministic handling
preserves exact data contents; identification itself may use an LLM.
Relevant locations receive symbolic labels, which name destinations
independently of their changing numerical offsets. An end-to-end model
that also regenerates data would be a different, unevaluated design.

The model returns a complete listing. Zone assembles it, restores retained
data, and checks the construction. The original binary, prepared listing,
and final binary are therefore different representations. The internal
LLM similarity score compares prepared and returned listings, excluding
changes already made during preparation.

The requested model is \texttt{deepseek-flash}; local version names come
from the client's alias table, without independent provider confirmation.
Both LLM roles request temperature 0.2, with no sampling seed or explicit
top-p. Repeated requests are not treated as deterministic reproductions.

\subsection{Prompting Conditions and Repeated Attempts}
A \emph{zero-shot} request describes the transformation and its constraints
without a worked example. A \emph{two-shot} request adds two examples of
input and transformed output. The number of shots thus counts examples
inside a prompt; it does not count generated variants or execution tests.
Each input, transformation family, and shot setting defines an experimental
\emph{cell}.

Validation feedback returns a rejected answer and its diagnostic---for
example, a missing label or assembly failure---and requests a correction.
It contains neither a functional success signal nor a scanner verdict.
The controlled benchmark disables correction and makes three independent
single-generation requests per cell, retaining failures. The original-corpus
campaign targets one output per cell, with some retries and failed
restarts in its logs. Its results are cell outcomes rather than uniform
first-generation success rates (Section~\ref{sec:methodology}).

\subsection{From a Proposed Listing to an Assessed Output}
Validation checks whether the response can be assembled and decoded,
whether tracked control-flow targets remain consistent, whether preserved
data survives, and whether the requested transformation's structural
requirements are met. A response can fail at any of these stages, or be
truncated before it contains a complete listing. Passing them means that
Zone accepts a constructed output, which then proceeds to the execution test.

An external \emph{test harness} launches the output and records its result.
It compares the complete expected buffer for controlled programs and
checks the class-specific observable for original payloads
(Section~\ref{sec:methodology}).


The functional test is applied after generation in these experiments. Saved records distinguish model requests, accepted outputs,
execution results, time, and estimated cost, allowing rejection and execution
failure to be analyzed independently.
